%
% Summary
% -------
\section{Summary}
\label{sec-summary}

The 1D nozzle solver provides a flexible Python interface for testing quantum linear equation
solvers on fluid flow type problems.
The core CFD solver, written in C, is linked via a library to enable developers to focus on their
quantum solver.
The benefit of the nozzle system is that it allows researchers to investigate the feedback between the 
inner linear system, solved using a quantum algorithm, and the outer classical non-linear update.

Although the solver is only 1 dimensional, there are a range of solver options that can be
used to generate matrices with up to 8 diagonals using the coupled compressible solver. The standard
pressure correction solver has 3 diagonals. The coupled incompressible solver has 6 diagonals and
gives rapid convergence which is useful when emulating the quantum circuit is time-consuming.

The plotting routines enable easy comparison of classical and quantum results.
